% Part: model-theory
% Chapter: interpolation
% Section: interpolation-proof

\documentclass[../../../include/open-logic-section]{subfiles}

\begin{document}

\olfileid[tr]{mod}{int}{prf}

\olsection{Craig Aradeğerleme Teoremi}

\begin{thm}[Craig aradeğerleme teoremi]
\ollabel{thm:interpol}
$\Entails !A\lif !B$ ise $\Entails !A\lif !C$ ve
$\Entails !C\lif !B$ olacak bir !!a{sentence}~$!C$ vardır; ayrıca
$!C$ içindeki her !!{constant}, !!{function} ve !!{predicate}
($\eq$ dışında) hem $!A$ hem $!B$ içinde geçer. !!{sentence}~$!C$
ifadesine $!A$ ile $!B$'nin bir \emph{aradeğerleyeni} denir.
\end{thm}

\begin{proof}
$\Lang L_1$, $!A$'nın; $\Lang L_2$, $!B$'nin dili olsun ve
$\Lang L_0=\Lang L_1\cap\Lang L_2$ alalım. Her
$i\in\{0,1,2\}$ için $\Lang L'_i$, $\Lang L_i$ diline sonsuz sayıda
yeni !!{constant}s~$\Obj c_0,\Obj c_1,\Obj c_2,\dots$ eklenerek elde
edilsin.

$!A$ sağlanamazsa $\lexists[x][\eq/[x][x]]$ bir aradeğerleyendir.
$\lnot !B$ sağlanamazsa (dolayısıyla $!B$ geçerliyse)
$\lexists[x][\eq[x][x]]$ bir aradeğerleyendir. Bu yüzden $!A$ ile
$\lnot !B$'nin her ikisinin de sağlanabilir olduğunu varsayabiliriz.

Aradeğerleme teoreminin karşıt tersini ispatlamak için $!A$ ile $!B$
için hiçbir aradeğerleyen bulunmadığını varsayalım. Başka bir deyişle
$\{!A\}$ ile $\{\lnot !B\}$ kümeleri $\Lang L_0$ içinde ayrılamaz
olsun.

Amacımız $(\{!A\},\{\lnot !B\})$ çiftini maksimal ayrılamaz bir
$(\Gamma^*,\Delta^*)$ çiftine genişletmektir. $!A_0,!A_1,!A_2,\dots$,
$\Lang L_1$ dilinin; $!B_0,!B_1,!B_2,\dots$ ise $\Lang L_2$ dilinin
!!{sentence}s ifadelerini numaralandırsın. $n\geq0$ için artan iki
!!{sentence}s kümeleri dizisi $(\Gamma_n,\Delta_n)$ tanımlayalım.
$\Gamma_0=\{!A\}$ ve $\Delta_0=\{\lnot !B\}$ alınsın.
$(\Gamma_n,\Delta_n)$ tanımlandıysa $\Gamma_{n+1}$ ve
$\Delta_{n+1}$ şöyle tanımlansın:
\begin{enumerate}
\item $\Gamma_n\cup\{!A_n\}$ ile $\Delta_n$, $\Lang L'_0$ içinde
  ayrılamazsa $!A_n$'yi $\Gamma_{n+1}$ içine koyun. Ayrıca $!A_n$,
  varoluşsal !!{formula}~$\lexists[x][!S]$ ise $\Gamma_n$,
  $\Delta_n$, $!A_n$ veya $!B_n$ içinde geçmeyen yeni bir
  !!{constant}~$c$ seçip $\Subst{!S}{c}{x}$ ifadesini
  $\Gamma_{n+1}$ içine koyun.
\item $\Gamma_{n+1}$ ile $\Delta_n\cup\{!B_n\}$, $\Lang L'_0$
  içinde ayrılamazsa $!B_n$'yi $\Delta_{n+1}$ içine koyun. Ayrıca
  $!B_n$, varoluşsal !!{formula}~$\lexists[x][!S]$ ise
  $\Gamma_{n+1}$, $\Delta_n$, $!A_n$ veya $!B_n$ içinde geçmeyen yeni
  bir !!{constant}~$c$ seçip $\Subst{!S}{c}{x}$ ifadesini
  $\Delta_{n+1}$ içine koyun.
\end{enumerate}
Son olarak
\begin{align*}
  \Gamma^* & = \bigcup_{n\ge 0} \Gamma_n, &
  \Delta^* & = \bigcup_{n\ge 0} \Delta_n
\end{align*}
tanımlayalım. Şimdi $n$ üzerinde eşzamanlı tümevarımla şunları
ispatlayabiliriz:
\begin{enumerate}
\item\ollabel{part-a} $\Gamma_n$ ile $\Delta_n$, $\Lang L'_0$ içinde
  ayrılamazdır;
\item\ollabel{part-b} $\Gamma_{n+1}$ ile $\Delta_n$, $\Lang L'_0$
  içinde ayrılamazdır.
\end{enumerate}
\olref{part-a} için temel durum \olref[sep]{lem:sep1} ile verilir.
\olref{part-b} için üç durumu ayırmalıyız:
\begin{enumerate}
\item $\Gamma_0\cup\{!A_0\}$ ile $\Delta_0$ ayrılabilirse
  $\Gamma_1=\Gamma_0$ ve \olref{part-b}, \olref{part-a} ile aynıdır;
\item $\Gamma_1=\Gamma_0\cup\{!A_0\}$ ise kuruluş gereği
  $\Gamma_1$ ile $\Delta_0$ ayrılamazdır;
\item $!A_0$ varoluşsal ve
  $\Gamma_1=\Gamma_0\cup\{\lexists[x][!S],\Subst{!S}{c}{x}\}$ ise,
  kuruluş gereği $\Gamma_0\cup\{\lexists[x][!S]\}$ ile $\Delta_0$
  ayrılamazdır. \olref[sep]{lem:sep2} gereği
  $\Gamma_0\cup\{\lexists[x][!S],\Subst{!S}{c}{x}\}$ ile
  $\Delta_0$ da ayrılamazdır.
\end{enumerate}
Bu, yukarıdaki \olref{part-a} ve \olref{part-b} için tümevarım temelini
tamamlar. Şimdi tümevarım adımına geçelim. \olref{part-a} için
$\Delta_{n+1}=\Delta_n\cup\{!B_n\}$ ise $\Gamma_{n+1}$ ile
$\Delta_{n+1}$ kuruluş gereği ayrılamazdır ($!B_n$ varoluşsal olsa
bile \olref[sep]{lem:sep2} gereği). $\Delta_{n+1}=\Delta_n$ ise
($\Gamma_{n+1}$ ile $\Delta_n\cup\{!B_n\}$ ayrılabildiği için)
\olref{part-b} tümevarım varsayımını kullanırız. \olref{part-b} için
$\Gamma_{n+2}=\Gamma_{n+1}\cup\{!A_{n+1}\}$ ise $\Gamma_{n+2}$ ile
$\Delta_{n+1}$ kuruluş gereği ayrılamazdır ($!A_{n+1}$ varoluşsal
olsa bile \olref[sep]{lem:sep2} gereği); $\Gamma_{n+2}=\Gamma_{n+1}$
ise az önce ispatlanan \olref{part-a} tümevarım adımını kullanırız.
Böylece \olref{part-a} ve \olref{part-b} üzerindeki tümevarım biter.

Buradan $\Gamma^*$ ile $\Delta^*$'nın ayrılamaz olduğu çıkar. Aksi
halde bir ayırma tümcesi için gereken sonlu öncüller bazı
$\Gamma_n,\Delta_n$ içinde bulunur ve \olref{part-a} ile çelişirdi.
Özellikle $\Gamma^*$ ile $\Delta^*$ tutarlıdır: Birincisi veya ikincisi
tutarsız olsaydı sırasıyla $\lexists[x][\eq/[x][x]]$ veya
$\lforall[x][\eq[x][x]]$ tarafından ayrılırlardı.

Şimdi $\Gamma^*$'ın $\Lang L'_1$ içinde, benzer biçimde $\Delta^*$'ın
$\Lang L'_2$ içinde maksimal tutarlı olduğunu gösterelim. Birincisi
için bazı $n\geq0$ bakımından $!A_n\notin\Gamma^*$ ve
$\lnot !A_n\notin\Gamma^*$ varsayalım. $!A_n\notin\Gamma^*$ ise
$\Gamma_n\cup\{!A_n\}$ ile $\Delta_n$ ayrılabilir ve
$!C\in\Lang L'_0$ olacak biçimde
\begin{align*}
  \Gamma^* & \Entails !A_n \lif !C, &
  \Delta^* & \Entails \lnot !C
\end{align*}
olur. Benzer biçimde $\lnot !A_n\notin\Gamma^*$ ise
$!C'\in\Lang L'_0$ vardır ve
\begin{align*}
  \Gamma^* & \Entails \lnot !A_n \lif !C', &
  \Delta^* & \Entails \lnot !C'.
\end{align*}
Önermeler mantığı gereği $\Gamma^*\Entails !C\lor !C'$ ve
$\Delta^*\Entails\lnot(!C\lor !C')$; dolayısıyla $!C\lor !C'$,
$\Gamma^*$ ile $\Delta^*$'yı ayırır. Benzer bir sav $\Delta^*$'ın
maksimal olduğunu gösterir.

Son olarak $\Gamma^*\cap\Delta^*$'nın $\Lang L'_0$ içinde maksimal
tutarlı olduğunu gösterelim. Tutarlı kümelerin kesişimi olduğundan
tutarlı olduğu açıktır. Maksimalliği göstermek için
$!S\in\Lang L'_0$ alın. $\Gamma^*$,
$\Lang{L'_1}\supseteq\Lang{L'_0}$ içinde; $\Delta^*$ da
$\Lang{L'_2}\supseteq\Lang{L'_0}$ içinde maksimaldir. Bu nedenle
$!S\in\Gamma^*$ ya da $\lnot !S\in\Gamma^*$ ve
$!S\in\Delta^*$ ya da $\lnot !S\in\Delta^*$. Eğer
$!S\in\Gamma^*$ ve $\lnot !S\in\Delta^*$ ise $!S$ bu iki kümeyi
ayırır; $\lnot !S\in\Gamma^*$ ve $!S\in\Delta^*$ ise $\lnot !S$
ayırır. Dolayısıyla ya $!S\in\Gamma^*\cap\Delta^*$ ya da
$\lnot !S\in\Gamma^*\cap\Delta^*$; böylece kesişim maksimaldir.

$\Gamma^*$ maksimal tutarlı olduğundan, tamlık teoreminin ispatındaki
gibi (\olref[fol][com][cth]{thm:completeness}), !!{domain}~$\Domain{M'_1}$
tam ve yalnızca !!{constant}s yorumlayan $\Assign c{M'_1}$
elemanlarından oluşan bir $\Struct{M'_1}$ modeli vardır. Benzer biçimde
$\Delta^*$'ın, !!{domain}~$\Domain{M'_2}$ sabit yorumları
$\Assign c{M'_2}$ ile verilen bir $\Struct{M'_2}$ modeli vardır.

$\Struct{M_1}$, $\Struct{M'_1}$ yapısından
$\Lang L'_1\setminus\Lang L'_0$ içindeki !!{constant}s, !!{function}s
ve !!{predicate}s yorumları atılarak; $\Struct{M_2}$ de benzer biçimde
elde edilsin. O zaman
$h(\Assign c{M'_1})=\Assign c{M'_2}$ ile tanımlanan
$h\colon M_1\to M_2$, $\Lang L'_0$ içinde bir izomorfizmadır; çünkü
gösterildiği üzere $\Gamma^*\cap\Delta^*$ bu dilde maksimal tutarlıdır.
Gerçekten her $\Lang L'_0$-!!{sentence} ya hem $\Gamma^*$ hem
$\Delta^*$ içinde bulunur ya da hiçbirinde bulunmaz. Dolayısıyla
$\Assign c{M'_1}\in\Assign P{M'_1}$ ancak ve ancak
$\Atom P c\in\Gamma^*$, ancak ve ancak $\Atom P c\in\Delta^*$, ancak
ve ancak $\Assign c{M'_2}\in\Assign P{M'_2}$. İzomorfizmanın öteki
koşulları da benzer biçimde gösterilir.

Şimdi !!{language}~$\Lang L_1\cup\Lang L_2$ için bir $\Struct M$
modelini şöyle tanımlayalım:
\begin{enumerate}
\item !!{domain}~$\Domain M$, tam olarak $\Domain{M_2}$, yani bütün
  $\Assign c{M'_2}$ elemanlarının kümesidir;
\item !!a{predicate}~$P$, $\Lang L_2\setminus\Lang L_1$ içindeyse
  $\Assign P M=\Assign P{M'_2}$;
\item bir $P$ yüklemi $\Lang L_1\setminus\Lang L_2$ içindeyse
  $\Assign P M=h(\Assign P{M'_1})$; yani
  $\tuple{\Assign{c_1}{M'_2},\dots,\Assign{c_n}{M'_2}}\in\Assign P M$
  ancak ve ancak
  $\tuple{\Assign{c_1}{M'_1},\dots,\Assign{c_n}{M'_1}}\in\Assign P{M'_1}$;
\item !!a{predicate}~$P$, $\Lang L_0$ içindeyse
  $\Assign P M=\Assign P{M'_2}=h(\Assign P{M'_1})$;
\item $\Lang L_1\cup\Lang L_2$ dilinin !!^{function}s, !!{constant}s
  dâhil, benzer biçimde ele alınır.
\end{enumerate}

Son olarak !!{formula}s üzerinde tümevarımla $\Struct M$ yapısının
$\Lang L'_1$ dilindeki bütün !!{formula}s bakımından $\Struct{M'_1}$
ile, $\Lang L'_2$ dilindekiler bakımından $\Struct{M'_2}$ ile
uyuştuğu gösterilir. Özellikle
$\Sat{M}{\Gamma^*\cup\Delta^*}$; buradan $\Sat{M}{!A}$ ve
$\Sat{M}{\lnot !B}$, dolayısıyla $\not\Entails !A\lif !B$ çıkar.
Bu, Craig aradeğerleme teoreminin ispatını tamamlar.
\end{proof}

\end{document}
